--- docs/design/proposal/manuscript/sections/standards-methods.tex

+++ docs/design/proposal/manuscript-intragroup-fx-20261013/sections/standards-methods.tex

@@ -2,5 +2,5 @@

 \section{Rule selection}\label{method-standards-methods-rule-selection}}
 
-The proposed selection procedure distinguishes operative authority from historical source support. The retained sources include official AASB compilations, the original EU incorporating regulation and an official ASBE 21 publication \citep{aasb9dated,aasb15dated,aasb16dated,aasb112dated,eu20231803,mofasbe21}. Historical US lease additions are separately identified in the retained amendment source \citep{fasbasu201602}.
+The proposed selection procedure distinguishes operative authority from historical source support. The retained sources include official AASB compilations, the original EU incorporating regulation and official ASBE lease and income-tax texts \citep{aasb9dated,aasb15dated,aasb16dated,aasb112dated,eu20231803,mofasbe21,moflease2018,mofasbe18,mofinterpretation16}. Historical US lease additions are separately identified in the retained amendment source \citep{fasbasu201602}.
 
 AuthorityRequirement contains issuer/framework, jurisdiction, entity class, annual-period/application interval, paragraph and local source Ref, incorporating instrument, amendment chain, transition, elections and unresolved dependencies. RuleSelection binds those requirements to the accounting event and evidence view. The algorithm evaluates chronological applicability first, then scope and facts, then available elections/transition, then measurement/presentation. If two supported rules can apply, retain the explicit alternative and required election; if one is contradicted by facts, mark ineligible. Missing current authority gives unresolved, never eligibility through a familiar rule name.
@@ -29,13 +29,151 @@

 
 \hypertarget{asbe-official-lease-support-and-remaining-tax-sources}{%
-\section{ASBE official lease support and remaining tax sources}\label{method-standards-methods-asbe-official-lease-support-and-remaining-tax-sources}}
-
-The official Ministry of Finance ASBE 21 page was successfully retrieved in this source investigation. The retained HTML and extracted text preserve the original Chinese wording. Page publication is 28 October 2019; article 68 states effect from 1 January 2019. The separate 2018 issuance notice with entity-specific staged application remains inaccessible, so the page alone does not establish every entity's adoption date or the complete amendment chain. The saved preparation packet's secondary reproduction remains historical evidence, with this new official source taking precedence for the actual inspected wording.
-
-Articles 4--8 require control of identified use, economic benefits and direction; article 6 includes both conditions for substantive substitution and the express presumption when difficult to determine. Articles 14--18 specify commencement, ROU cost and PV, using implicit rate or incremental borrowing rate if the former cannot be determined. Articles 20--23 separate cost-model depreciation and periodic interest; articles 25--29 distinguish reassessment/index changes/modifications and their discount-rate rules. Articles 30--33 define the exemptions; a two-year term is not short-term. Articles 60--67 specify transition choices and disclosures; the chosen transition cannot be inferred from a present closing liability.
+\section{ASBE lease adoption and deferred tax}\label{method-standards-methods-asbe-official-lease-support-and-remaining-tax-sources}}
+
+For an ASBE lease, selecting the lease standard and selecting the related income-tax treatment require different application dates. Article 68 of the revised ASBE 21 states a general date of 1 January 2019, but the Ministry of Finance's 2018 issuance notice specifies the enterprise groups. Paragraph I requires application from that date for enterprises listed both domestically and overseas, and for overseas-listed enterprises preparing financial statements under IFRS or ASBE. Other enterprises applying ASBE must apply the revision from 1 January 2021 \citep[issuance notice, para.\ I]{moflease2018}.
+
+Paragraph II permits earlier application for enterprises whose parent or subsidiary is listed overseas and prepares its overseas financial statements under IFRS or ASBE. Their early date cannot precede simultaneous application of the specified 2017 financial-instruments and revenue standards. Paragraph III discontinues the 2006 lease standard and its application guidance when the enterprise applies the revision. Thus a 2020 reporting date alone does not establish use of the revised lease standard for an enterprise in the second group: its qualifying early adoption would need evidence. The retained official Chinese texts include the 2018 notice and attachment and the Ministry's later ASBE 21 page, published on 28 October 2019 \citep{moflease2018,mofasbe21}. They replace the earlier secondary reproduction as support for the inspected wording.
+
+Articles 4--8 require control of identified use, economic benefits and direction; article 6 includes both conditions for substantive substitution and the express presumption when difficult to determine. Articles 14--18 specify commencement, ROU cost and PV, using implicit rate or incremental borrowing rate if the former cannot be determined. Articles 20--23 separate cost-model depreciation and periodic interest; articles 25--29 distinguish reassessment/index changes/modifications and their discount-rate rules. Articles 30--33 define the exemptions; a two-year term is not short-term. Articles 60--67 specify transition choices and disclosures; the chosen transition cannot be inferred from a present closing liability \citep[attached standard, arts.\ 4--8, 14--33, 60--67]{moflease2018}.
 
 For the declared base lease, these inspected recognition/measurement provisions support proposed ROU 210/105/0 and liability 210/110/0 with expense 126/116, under the same stated facts, depreciation assumption and fixed rule version. A same-asset modification is tested against 28--29, rather than labeled a separate lease without increased scope and commensurate consideration. A 2019 first-adopter using article 61(2) records cumulative opening changes without adjusting comparatives, while 61(1) is retrospective; comparison datasets must retain that distinction and the entity's actual election.
 
-Official ASBE 18 and Interpretation 16 pages still returned 502. The retained official 2023 MoF tax application case explicitly discusses equal lease temporary differences and Interpretation 16; it supplies a precise proposed answer for this case, not a complete current income-tax standard or real tax law. Reproducing the old initial-recognition exception from ASBE 18 without the later interaction would contradict this official application evidence. Keep gross lease DTA/DTL and the admitted tax-base/recovery assumptions in proposed calculations, but require operative text, effective-period/transition and actual law before real permission. Full ASBE 22 credit, consolidated statements, group combinations, and insurance revisions likewise retain named source dependencies; do not assert full IFRS equivalence from similar amounts.
+The tax question is why equal opening lease balances can still require separate deferred-tax assets and liabilities. The Ministry's original ASBE 18 defines tax bases and temporary differences in articles 4--9 and contains initial-recognition exceptions in articles 11 and 13. Article 13 also limits a deferred-tax asset to the probable taxable profit available to use the deduction. Articles 16--23 govern measurement, recovery or settlement assumptions, subsequent review, allocation and noncurrent classification \citep{mofasbe18}. The original text must be read with later interpretations.
+
+Interpretation 16 removes the initial-recognition exemption for a single transaction that is not a business combination, initially affects neither accounting profit nor taxable income or loss, and creates equal taxable and deductible temporary differences. Section I(a) expressly includes qualifying initial lease liabilities and right-of-use assets; the corresponding deferred-tax liabilities and assets are recognized separately under ASBE 18. Equal carrying amounts alone do not establish equal temporary differences, and the asset's recoverability condition still applies. Section IV makes this treatment effective from 1 January 2023; Section I(b) permits early application from the year of issuance, 2022, with disclosure \citep[secs.\ I(a)--(b), IV]{mofinterpretation16}. This tax date is distinct from the enterprise's lease-standard adoption date. The operative text now supports the treatment previously illustrated by the retained official 2023 tax application case.
+
+For the synthetic two-year lease, retain the assumed payments of 121, 10\% discount rate and straight-line depreciation. All amounts use the same illustrative monetary unit. Assume that tax deductions are allocated to the liability, both tax bases are zero, the applicable tax rate is 25\%, and sufficient taxable profit satisfies the deferred-tax asset recognition condition. Also assume that the transaction meets Section I(a)'s conditions and falls within its application period. At commencement the ROU asset and liability are each 210. The taxable ROU difference gives a deferred-tax liability (DTL) of 52.5; the deductible liability difference gives a deferred-tax asset (DTA) of 52.5. The entry debits DTA and credits DTL by those amounts. At the first annual close, liability 110 and ROU 105 give DTA 27.5 and DTL 26.25. Reducing the DTA by 25 creates deferred-tax expense; reducing the DTL by 26.25 creates a benefit, leaving net benefit 1.25. Both balances reverse by expiry. Applying the old exemption to omit both opening deferred balances is wrong for these qualifying facts and dates: it conceals the required gross balances and the later tax effect.
+
+Transition also needs more than an opening journal entry. Section I(b) requires adjustments for qualifying transactions between the earliest comparative opening and application. For qualifying existing balances at that opening, cumulative effects adjust opening retained earnings and the relevant financial-statement items. Consider a separate transition example with cash 500, ROU 105, lease liability 110 and equity 495, with no deferred tax yet recognized. Under the same assumed tax bases, rate and recoverability, recognition debits DTA 27.5, credits DTL 26.25 and credits opening retained earnings 1.25. The resulting balance sheet is
+\[
+\underbrace{500+105+27.5}_{\text{assets }632.5}
+=\underbrace{110+26.25}_{\text{liabilities }136.25}
++\underbrace{495+1.25}_{\text{equity }496.25}.
+\]
+The equity adjustment records the cumulative effect at the comparative opening. If the remaining lease year falls within the comparative year, the DTA of 27.5 and DTL of 26.25 reverse during that year, giving comparative deferred-tax expense 1.25. The opening adjustment alone therefore does not complete the transition. Neither ASBE 21's article 61(2) lease-adoption election nor Interpretation 16 Section III's no-comparative-adjustment rule for its separate share-payment topic may be substituted for Section I(b)'s tax transition.
+
+These are gross recognition and measurement amounts. Separate recognition and noncurrent classification do not determine whether an actual balance sheet may offset DTA and DTL. That decision needs the applicable presentation provisions and entity/tax-authority facts. Actual tax law, tax bases and rates, recoverable profit, the enterprise's adoption and elections, and later amendments remain to be established before applying the example to an issuer. Full ASBE 22 credit, consolidated statements, group combinations and insurance revisions also retain specific source dependencies; similar amounts do not establish full IFRS equivalence.
+
+\clearpage
+\section{Reading an actual lease transition}
+\label{sec:pingan-lease-transition}
+
+Ping An's 2019 interim report shows what happens when the dated lease rules
+meet an actual disclosure. The group reports adoption of revised ASBE~21 on
+1 January 2019 and leaves its 2018 comparative financial statements
+unadjusted. Its stated listings in Shanghai and Hong Kong place it within
+the enterprise class covered by the 2018 notice's 2019 application date
+\citep[pp.\ 112--113]{pingan2019asbe};
+\citep[notice, para.\ I]{moflease2018}. This supplies an actual adoption
+fact for the representation. The later deferred-tax treatment described
+above still requires its own applicable date and tax facts.
+
+The immediate reconstruction problem is to keep three observations apart:
+the original December presentation, the January adoption balances and the
+June closing balances. A right-of-use asset records the group's recognised
+right to use leased property; the lease liability measures the associated
+recognised payment obligation. Table~\ref{tab:pingan-lease-balances} retains
+the amounts and their dates. All amounts in this example are displayed
+millions of renminbi; arithmetic on them does not establish exact underlying
+amounts or the issuer's rounding rule.
+
+\begin{table}[!ht]
+\centering
+\caption{Adoption balances and later carrying balances answer different questions.}
+\label{tab:pingan-lease-balances}
+\begin{tabular}{@{}lrr@{}}
+\toprule
+Observation & Right-of-use assets & Lease liabilities \\
+\midrule
+31 December 2018, unadjusted comparative & --- & --- \\
+1 January 2019, reported at adoption & 18,500 & 17,667 \\
+30 June 2019, carrying balance & 17,390 & 17,212 \\
+\bottomrule
+\end{tabular}
+\end{table}
+
+The comparative dashes reproduce the presentation under the previous basis;
+they do not say that the group had no lease contracts in December. The
+January amounts come from the policy-change note, whereas the June amounts
+come from the consolidated balance sheet
+\citep[pp.\ 99--100, 112]{pingan2019asbe}. Subtracting a December dash
+treated as zero from the June liability would mix recognition at transition
+with six months of subsequent activity.
+
+The opening asset exceeds the liability by $18{,}500-17{,}667=833$. That
+difference alone does not identify an increase in retained earnings.
+Article~61(2) permits specified alternative measurements of the right-of-use
+asset for former operating leases, including a liability-based amount
+adjusted for prepaid rent; its cumulative approach adjusts retained earnings
+and other relevant financial-statement items
+\citep[attachment, art.\ 61(2)]{moflease2018}. To see why the other items
+matter, suppose solely for illustration that pre-existing prepaid rent of
+833 is reclassified. Debiting the new right-of-use asset by 18,500 and
+crediting the new liability by 17,667 and prepaid rent by 833 balances the
+entry with no equity movement. Assigning the same residual entirely to
+retained earnings would instead require a different assumption about all
+other transition adjustments. The report does not identify the illustrative
+prepayment or supply a complete opening journal. Reconstruction therefore
+retains the unexplained components rather than creating an equity plug.
+
+The policy note nevertheless narrows the permitted interpretation. It says
+the opening liability discounts remaining payments at the lessee's
+incremental borrowing rate on 1 January 2019. It names five simplifications:
+treating operating leases with a remaining term shorter than twelve months
+as short-term leases; excluding initial direct costs at transition; using
+one discount rate for leases with similar characteristics; using actual
+option exercise and other recent information to determine the lease term;
+and assessing pre-existing onerous contracts instead of a separate
+right-of-use impairment test. These correspond to the specified provisions
+of Article~63, with the rate and direct-cost choices appearing in the same
+numbered provision \citep[p.\ 113]{pingan2019asbe};
+\citep[attachment, art.\ 63]{moflease2018}. The record should preserve
+the reported elections without assigning an unreported rate or deciding
+which asset-measurement alternative was used for each lease. The note also
+calls the difference between discounted prior operating-lease commitments
+and the recognised opening liability ``not material''. That qualification
+supplies neither an equality nor a numeric tolerance.
+
+The June maturity disclosure adds a different measurement of the obligation.
+Under its heading for undiscounted contractual and expected cash flows,
+the lease-liability row gives 1,515 within three months, 5,885 in three to
+twelve months, 11,048 in one to five years and 485 beyond five years. Those
+numeric cells sum to 18,933; the immediate and no-term columns contain
+dashes \citep[pp.\ 164--165]{pingan2019asbe}. The first two dated buckets
+sum to 7,400, but disclose no payment dates within those buckets. The
+18,933 total and the 17,212 carrying liability can coexist because they
+measure different things. Their difference, 1,721, is not the interest
+expense for the half-year and does not by itself determine a discount rate.
+
+Even on the post-adoption basis, a change in a balance is not a complete
+account of the transactions that produced it. The right-of-use asset falls
+by $17{,}390-18{,}500=-1{,}110$ and the liability by
+$17{,}212-17{,}667=-455$. Additions, remeasurements, disposals and other
+movements can accompany depreciation or payments. The first difference
+therefore does not identify depreciation, and the second does not identify
+cash paid. The consolidated cash-flow statement's finance-leasing business
+borrowing and receivable lines concern the group's financing activities as
+a business; their labels do not make them payments on its own lessee
+obligations \citep[pp.\ 106--107]{pingan2019asbe}.
+
+The same care prevents a false narrative conflict. In its major-contract
+discussion, the company says that it has no reportable custody, contracting,
+lease or other major-contract matters for the period
+\citep[p.\ 86]{pingan2019asbe}. Preserving the company reference and
+the reportable-major-contract qualification gives a different assertion
+from saying that the consolidated group has no leases. Removing those
+qualifiers would create the apparent contradiction with the balance sheet.
+
+For the proposed representation, the useful result is a set of dated,
+source-supported balances, policy elections and maturity restrictions.
+A query for the reported June carrying liability can return 17,212 with
+its basis and precision. A query for future scheduled payments must keep
+the disclosed maturity buckets and leave their internal dates unresolved.
+A query for the complete transition journal or the half-year's lessee cash
+payments must retain the missing inputs. These distinct answers let actual
+disclosures constrain histories without pretending that every balanced
+completion is the company's history.
+
+\clearpage
 
 \hypertarget{matched-case-outcomes-and-decision}{%
@@ -84,3 +222,3 @@

 \end{longtable}
 
-These decisions establish a common procedure and specific proposed treatments. The broader promised authority coverage remains incomplete. The new official ASBE lease support narrows that blocker; it does not erase the specific US, ASBE tax/sector, legal and current-applicability work. Feasible internal preparation continues without a paid engagement.
+These decisions establish a common procedure and specific proposed treatments. The recovered ASBE lease notice and tax interpretation support the stated application schedules and conditional example. Broader authority coverage remains incomplete: operative US sources, ASBE sector methods, actual tax law and presentation conditions, and the applicable amendment and entity-election evidence still require internal investigation.
